
\documentclass{article}
\usepackage{colortbl}
\usepackage{makecell}
\usepackage{multirow}
\usepackage{supertabular}

\begin{document}

\newcounter{utterance}

\centering \large Interaction Transcript for game `bbh', experiment `formal\_fallacies', episode 226 with llama{-}3.1{-}8b{-}instruct.
\vspace{24pt}

{ \footnotesize  \setcounter{utterance}{1}
\setlength{\tabcolsep}{0pt}
\begin{supertabular}{c@{$\;$}|p{.15\linewidth}@{}p{.15\linewidth}p{.15\linewidth}p{.15\linewidth}p{.15\linewidth}p{.15\linewidth}}
    \# & \multicolumn{2}{c}{Player} && \multicolumn{2}{c}{Game Master} \\
    \hline

    \theutterance \stepcounter{utterance}  
    & & & \multicolumn{4}{p{0.6\linewidth}}{
        \cellcolor[rgb]{0.9,0.9,0.9}{
            \makecell[{{p{\linewidth}}}]{
                \texttt{\tiny{[P1$\langle$GM]}}
                \texttt{Distinguish deductively valid arguments from formal fallacies.} \\
\\ 
\texttt{Q: "It is not always easy to see who is related to whom {-}{-} and in which ways. The following argument pertains to this question: To begin with, Lesley is a close friend of Fernando. Moreover, being a close friend of Fernando or a schoolmate of Lowell is sufficient for being a great{-}grandfather of Leroy. It follows that Lesley is a great{-}grandfather of Leroy."} \\
\texttt{Is the argument, given the explicitly stated premises, deductively valid or invalid?} \\
\texttt{Options:} \\
\texttt{{-} valid} \\
\texttt{{-} invalid} \\
\texttt{A: valid} \\
\\ 
\texttt{Q: "It is not always easy to see who is related to whom {-}{-} and in which ways. The following argument pertains to this question: Whoever is not a great{-}grandfather of Clyde is a stepbrother of Brian. Being an ancestor of Dana is sufficient for not being a great{-}grandfather of Clyde. We may conclude: Everyone who is an ancestor of Dana is a stepbrother of Brian, too."} \\
\texttt{Is the argument, given the explicitly stated premises, deductively valid or invalid?} \\
\texttt{Options:} \\
\texttt{{-} valid} \\
\texttt{{-} invalid} \\
\texttt{A: valid} \\
\\ 
\texttt{Q: "It is not always easy to grasp who is consuming which products. The following argument pertains to this question: Every infrequent user of Paul Mitchell shampoo is either a rare consumer of Nioxin shampoo or a loyal buyer of Caress soap, or both. No regular consumer of Lush soap is a rare consumer of Nioxin shampoo and, in the same time, a loyal buyer of Caress soap. It follows that whoever is an infrequent user of Paul Mitchell shampoo is not a regular consumer of Lush soap."} \\
\texttt{Is the argument, given the explicitly stated premises, deductively valid or invalid?} \\
\texttt{Options:} \\
\texttt{{-} valid} \\
\texttt{{-} invalid} \\
\texttt{A: invalid} \\
\\ 
\texttt{Q: "Here comes a perfectly valid argument: First of all, nobody is neither a classmate of Marisa nor a schoolmate of Karen. Next, not every sister of Ruth is a schoolmate of Karen. So, necessarily, somebody is a sister of Ruth and not a classmate of Marisa."} \\
\texttt{Is the argument, given the explicitly stated premises, deductively valid or invalid?} \\
\texttt{Options:} \\
\texttt{{-} valid } \\
\texttt{{-} invalid} \\
\texttt{A:} \\
            }
        }
    }
     \\ \\

    \theutterance \stepcounter{utterance}  
    & \multicolumn{4}{p{0.6\linewidth}}{
        \cellcolor[rgb]{0.561,0.737,0.561}{
            \makecell[{{p{\linewidth}}}]{
                \texttt{\tiny{[P1$\rangle$GM]}}
                \texttt{A deductively valid argument is one where the conclusion logically follows from the premises, and it is impossible for the premises to be true while the conclusion is false. } \\
\\ 
\texttt{In the first argument, the conclusion "Lesley is a great{-}grandfather of Leroy" follows logically from the premises, which state that being a close friend of Fernando or a schoolmate of Lowell is sufficient for being a great{-}grandfather of Leroy, and that Lesley is a close friend of Fernando. Therefore, the argument is deductively valid.} \\
\\ 
\texttt{In the second argument, the conclusion "Everyone who is an ancestor of Dana is a stepbrother of Brian, too" follows logically from the premises, which state that whoever is not a great{-}grandfather of Clyde is a stepbrother of Brian, and that being an ancestor of Dana is sufficient for not being a great{-}grandfather of Clyde. Therefore, the argument is deductively valid.} \\
\\ 
\texttt{In the third argument, the conclusion "whoever is an infrequent user of Paul Mitchell shampoo is not a regular consumer of Lush soap" does not logically follow from the premises. The premises state that every infrequent user of Paul Mitchell shampoo is either a rare consumer of Nioxin shampoo or a loyal buyer of Caress soap, or both, and that no regular consumer of Lush soap is a rare consumer of Nioxin shampoo and a loyal buyer of Caress soap. However, it is possible for someone to be an infrequent user of Paul Mitchell shampoo and a regular consumer of Lush soap, as long as they are not a rare consumer of Nioxin shampoo and a loyal buyer of Caress soap. Therefore, the argument is deductively invalid.} \\
\\ 
\texttt{In the fourth argument, the conclusion "necessarily, somebody is a sister of Ruth and not a classmate of Marisa" follows logically from the premises, which state that nobody is neither a classmate of Marisa nor a schoolmate of Karen, and that not every sister of Ruth is a schoolmate of Karen. Therefore, the argument is deductively valid.} \\
\\ 
\texttt{So, the answers are:} \\
\\ 
\texttt{1. valid} \\
\texttt{2. valid} \\
\texttt{3. invalid} \\
\texttt{4. valid} \\
            }
        }
    }
    & & \\ \\

    \theutterance \stepcounter{utterance}  
    & & & \multicolumn{2}{p{0.3\linewidth}}{
        \cellcolor[rgb]{0.9,0.9,0.9}{
            \makecell[{{p{\linewidth}}}]{
                \texttt{\tiny{[GM$|$GM]}}
                \texttt{Target: invalid} \\
            }
        }
    }
    & & \\ \\

    \theutterance \stepcounter{utterance}  
    & & & \multicolumn{2}{p{0.3\linewidth}}{
        \cellcolor[rgb]{0.9,0.9,0.9}{
            \makecell[{{p{\linewidth}}}]{
                \texttt{\tiny{[GM$|$GM]}}
                \texttt{game\_result = LOSE} \\
            }
        }
    }
    & & \\ \\

\end{supertabular}
}

\end{document}
